% Ad Atticum 7.13 (ad-atticum-07-13) --- English translation
% Translated by Claude (Anthropic) from the Latin (Perseus).
% Style guide: STYLE.md.

\ciceroLetterOpener{CICERO ATTICO salutem}

\ciceroSection{1} On the Vennonian business I agree with you. Labienus I judge a \textit{[Greek: hero]}. No act of civil life has long been more brilliant: who, granting nothing else, has at least accomplished this --- he has given pain to that man. But I think something has been accomplished, too, towards the sum of things. I have affection for Piso, too. His judgement of his son-in-law I suspect will be felt as a heavy blow. Yet you see what kind of a war this is. So far civil that it has been born not of a quarrel between citizens but of the audacity of one ruined citizen. He, moreover, is strong in his army, holds many by hope and promises, and has reached out for everything of everyone. To him the City has been handed over, stripped of any garrison, packed with supplies. What is there you should not fear from the man who reckons those temples and roofs not as his fatherland but as his plunder? What he is about to do, or in what fashion, I do not know --- without a senate, without magistrates. He will not even be able to pretend at anything \textit{[Greek: in a statesmanlike way]}. As for us, where can we lift ourselves up, or when? Whose leader is how \textit{[Greek: unstrategic]} you too can see --- a man to whom not even Picenum was a known country; and how unplanned the whole thing is, the facts themselves bear witness. To pass over the failings of ten years, what bargain has not been better than this present flight?

\ciceroSection{2} And as for what he is thinking now, I do not know, nor do I cease to inquire of him by letter. Nothing, it is settled, is more timorous; nothing more disordered. Accordingly I see neither garrison --- for the levying of which he was kept near the City --- nor place and seat for any garrison. All his hope rests on two legions, almost not his own, treacherously kept back. For so far the levy is of unwilling men, men who shrink from a fight; and the time for terms has been lost. What is to come I do not see; what is certain is, the fault is ours --- or our leader's --- for having put out of harbour without rudders and given ourselves over to the storm.

\ciceroSection{3} So as for our two Ciceros, I am in doubt what to do; for at times it seems best to send them off to Greece. As for Tullia and Terentia, when the coming of barbarians to the City is set before me, I am all in fear; but when Dolabella comes to mind, I breathe a little. I should like you to think over what should be done --- first \textit{[Greek: regarding what is fitting in itself]} (for I must take counsel about them on different grounds than about myself), then with an eye to opinion, lest we be reproached for having wanted them at Rome in the general flight of decent men. Indeed, you and Peducaeus too (for he has written to me) must consider what you yourselves are to do. For yours is such standing that the same things are demanded of you as of the most eminent citizens. But you will see to this --- since it is about me and mine that I want your reflection.

\ciceroSection{4} It remains for you both to investigate, so far as you can, what is going on, and write to me; and what you yourself reach by conjecture --- which I look for from you even more. For I look to you for the future acts that everyone else is reporting. As for the \textit{[Greek: prophet]} --- you will forgive my talkativeness, which both eases me when writing to you and draws out your letters. The riddle of the Oppii from Velia I plainly did not understand: it is more obscure than the number of Plato.
